\section{Related work}

\subsection{Frame selection for corrective prompting}

Frame selection for efficient human annotation has been studied in both active learning and interactive video segmentation.
Active-learning approaches often prioritise frames using criteria such as uncertainty, representativeness, or diversity, with the acquired annotations subsequently used to improve the model \cite{vijayanarasimhan2012active, kuznetsova2021efficient, rana2022all}.
For example, HVSA \cite{hvsa} combines uncertainty- and diversity-based selection with video-specific fine-tuning.

Interactive video segmentation considers a more directly related setting, in which user corrections are used to improve the propagated masks of the current video while the segmentation model remains fixed \cite{yao2020video}.
Within this setting, several methods have studied where and when the user should intervene.
IVOS-W \cite{IVOSW} recommends the next frame according to the expected video-level improvement after annotation, showing that the frame with the poorest current mask is not necessarily the most useful one to correct.
EVA-VOS \cite{delatolas2024learning} selects the next frame using representations derived from the image and predicted mask, favouring frames that differ from those already annotated.
Lazy-XMem \cite{vujasinovic2024strike} instead monitors uncertainty in the propagated masks and triggers a correction when the prediction becomes unreliable.
Although these methods improve where the user intervenes, they select correction frames sequentially. When multiple corrections are needed for a video, this requires repeated cycles of propagation, expert review, and frame selection before all necessary corrections are identified \cite{delatolas2024learning}.
Selecting a set of correction frames for each video can reduce this back-and-forth process while allowing the number of requested frames to adapt to the current propagation quality and the cost of expert assistance.

\subsection{Learning to defer for promptable video segmentation}

Learning to defer framework decides whether to retain a machine prediction or request expert assistance. Conventional formulations learn this decision by balancing the expected errors of the machine and expert against the cost of expert consultation \cite{mao2023two, mozannar2020consistent,verma2022calibrated,hemmer2304learning}.

Related human–AI collaboration methods extend this idea to settings with multiple experts. LECODU \cite{zhang2024learning} predicts how many experts should be consulted for each input and then randomly samples that number from the available expert pool before combining their predictions with the machine output. Top-k \cite{montreuil2025ask} learning to defer instead ranks experts according to their estimated suitability for the input and selects the highest-ranked k. Its adaptive Top-k(x) formulation additionally allows the consultation cardinality to vary across inputs.

Most L2D methods target classification.
L2RP~\cite{rasanjalee2026understanding} is, to our knowledge, the first application of L2D to promptable video segmentation, but it defers at most once per video.
Extending the setting to multiple corrective prompts is not straightforward, because each correction re-propagates through the video and changes how much the remaining candidates would add.
This also prevents direct application of Top-\(k(x)\)~\cite{montreuil2025ask}, which scores each expert independently and can therefore select several adjacent frames that address the same failure.
Our framework instead selects the corrective frames jointly, conditioning each choice on the frames already selected, and allows the number of corrections to vary with the video.
